% pinn_slides.tex -- additional Physics-Informed Neural Network slides
% Source context: MIKU-7437/AI_Termoelastic_Waves_Geology, chapter 5 PINN file.

\begin{frame}{PINN: Physics-Informed Baseline}
  \begin{columns}[T]
    \begin{column}{0.48\textwidth}
      \begin{block}{Role in the comparison}
        \footnotesize
        The Physics-Informed Neural Network (PINN) is the route whose loss
        explicitly contains thermoelastic equation residuals.
      \end{block}
      \begin{itemize}
        \footnotesize
        \item Candidate predictor for the shared task.
        \item Physics-informed baseline for the data-driven routes.
        \item Uses supervised data and governing-equation constraints.
      \end{itemize}
    \end{column}
    \begin{column}{0.48\textwidth}
      \begin{alertblock}{Interpretation}
        \footnotesize
        PINN is not an exact numerical solver. It is a neural surrogate whose
        training objective encourages physically consistent predictions.
      \end{alertblock}
    \end{column}
  \end{columns}
\end{frame}

\begin{frame}{PINN Input--Output Mapping}
  \begin{block}{Pointwise thermoelastic prediction}
    \[
      f_{\theta}(x,y,t,E,\nu,\rho,\alpha,k,C_p)\rightarrow(T,u,v)
    \]
  \end{block}
  \begin{columns}[T]
    \begin{column}{0.48\textwidth}
      \footnotesize
      \begin{itemize}
        \item $(x,y)$: in-plane coordinate.
        \item $t$: observation time.
        \item $E,\nu,\rho,\alpha,k,C_p$: material parameters.
      \end{itemize}
    \end{column}
    \begin{column}{0.48\textwidth}
      \footnotesize
      \begin{itemize}
        \item $T$: predicted temperature.
        \item $u,v$: in-plane displacement components.
        \item Stress and strain are recovered from derivatives.
      \end{itemize}
    \end{column}
  \end{columns}
\end{frame}

\begin{frame}{PINN Network Structure}
  \begin{columns}[T]
    \begin{column}{0.50\textwidth}
      \begin{block}{Multilayer perceptron}
        \[
          h^{(\ell)}=\varphi\left(W^{(\ell)}h^{(\ell-1)}+b^{(\ell)}\right)
        \]
        \[
          h^{(0)}=(x,y,t,E,\nu,\rho,\alpha,k,C_p)
        \]
      \end{block}
    \end{column}
    \begin{column}{0.46\textwidth}
      \footnotesize
      The implementation uses a smooth activation such as $\tanh$, because the
      PDE residuals require derivatives of the network output with respect to
      space and time.
      \vskip 4pt
      \scriptsize
      $\theta$ collects all trainable weights and biases; automatic
      differentiation provides the derivatives used in the loss.
    \end{column}
  \end{columns}
\end{frame}

\begin{frame}{PINN Composite Loss}
  \begin{block}{Training objective}
    \[
      \mathcal{L}=\lambda_{\mathrm{sup}}\mathcal{L}_{\mathrm{sup}}
      +\lambda_{\mathrm{vel}}\mathcal{L}_{\mathrm{vel}}
      +\lambda_{\mathrm{el}}\mathcal{L}_{\mathrm{elastic}}
      +\lambda_{\mathrm{th}}\mathcal{L}_{\mathrm{thermal}}
    \]
  \end{block}
  \begin{columns}[T]
    \begin{column}{0.48\textwidth}
      \footnotesize
      \begin{block}{Data terms}
        \begin{itemize}
          \item supervised error for $(T,u,v)$;
          \item velocity consistency with exported COMSOL velocities.
        \end{itemize}
      \end{block}
    \end{column}
    \begin{column}{0.48\textwidth}
      \footnotesize
      \begin{block}{Physics terms}
        \begin{itemize}
          \item elastic-wave residual;
          \item coupled thermal residual.
        \end{itemize}
      \end{block}
    \end{column}
  \end{columns}
\end{frame}

\begin{frame}{PINN Thermoelastic Residuals}
  \begin{columns}[T]
    \begin{column}{0.52\textwidth}
      \begin{block}{Momentum balance}
        \[
          \rho\,\partial_t^2 u-\nabla\!\cdot\!\sigma=0
        \]
        \[
          \sigma=\lambda(\nabla\!\cdot\!u)I+2\mu\varepsilon-
          \gamma\theta I
        \]
      \end{block}
    \end{column}
    \begin{column}{0.44\textwidth}
      \begin{block}{Coupled heat equation}
        \[
          \rho C_p\partial_tT-k\nabla^2T+
          \gamma T_0\partial_t(\nabla\!\cdot\!u)=0
        \]
      \end{block}
      \scriptsize
      $\theta=T-T_{\mathrm{ref}}$ and $\gamma=3K\alpha$ connect temperature,
      strain, and stress inside the residual terms.
    \end{column}
  \end{columns}
\end{frame}

\begin{frame}{PINN Data Use and Resulting Role}
  \begin{columns}[T]
    \begin{column}{0.48\textwidth}
      \begin{block}{Training data}
        \footnotesize
        PINN uses the same COMSOL-derived 2D datasets: coordinates, time,
        material parameters, and ground-truth temperature and displacement.
      \end{block}
      \footnotesize
      Residuals are evaluated at collocation points using automatic
      differentiation.
    \end{column}
    \begin{column}{0.48\textwidth}
      \begin{block}{Why it matters}
        \footnotesize
        Because governing-equation residuals are part of training, PINN is the
        physically grounded reference route in the comparison framework.
      \end{block}
      \footnotesize
      Its outputs should still be interpreted within the simplified 2D
      assumptions and available training-data coverage.
    \end{column}
  \end{columns}
\end{frame}
